\section{DP-Quiz}\label{sec:13-dp-quiz}

\subsection{Welche Aussage ist korrekt?}\label{sec:13-welche-aussage-ist-korrekt}

\begin{contentbox}
Welche der folgenden Aussagen ist wahr?
\begin{factlist}
    \item[c.] Dynamische Programmierung eignet sich am besten für Probleme mit optimaler Substruktur und überlappenden Teilproblemen.
\end{factlist}
\end{contentbox}


\begin{contentbox}
Welche der folgenden Aussagen über dynamische Programmierung sind wahr?
\begin{factlist}
    \item[(I)] Dynamische Programmierung ist nur anwendbar, wenn man mit einem 2D-Array arbeitet oder eine 2D-DP-Tabelle aufbaut.
    \item[(II)] Dynamische Programmierung benötigt mindestens zwei verschachtelte for-Schleifen.
    \item d. Keine der Aussagen ist wahr.
\end{factlist}
\end{contentbox}


\begin{contentbox}
Welche der folgenden Aussagen über dynamische Programmierung sind wahr?
\begin{factlist}
    \item[(I)] Wenn ein Algorithmus eine zweidimensionale Matrix befüllt, dann handelt es sich um dynamische Programmierung.
    \item[(II)] Wenn ein Algorithmus verschachtelte for-Schleifen benutzt, dann ist es dynamische Programmierung.
    \item d. Keine der Aussagen ist wahr.
\end{factlist}
\end{contentbox}


\begin{contentbox}
Du und dein Kollege haben unabhängig voneinander DP-Lösungen für dasselbe Problem entwickelt. Beim Vergleich stellst du fest, dass dein Kollege andere Teilprobleme und eine DP-Tabelle mit mehr Dimensionen verwendet hat. Was beschreibt die Korrektheit eurer Lösungen am besten?
\begin{factlist}
    \item a. Es ist möglich, dass beide korrekte, valide und effiziente DP-Lösungen haben.
\end{factlist}
\end{contentbox}


\begin{contentbox}
Welche der folgenden Aussagen über dynamische Programmierung und Optimierungsprobleme sind immer wahr?
\begin{factlist}
    \item[(I)] Dynamische Programmierung befasst sich mit Optimierungsproblemen.
    \item[(II)] Wenn ein Problem ein Optimierungsproblem ist, kann dynamische Programmierung zur Lösung verwendet werden.
    \item b. Nur (I) ist wahr.
\end{factlist}
\end{contentbox}


\begin{contentbox}
Die $n$-te Fibonacci-Zahl $F(n)$ ist definiert als $F(n) = F(n-1) + F(n-2)$ mit den Basisfällen $F(1) = 1$ und $F(2) = 2$. Betrachte einen \textbf{Top-down}-DP-Algorithmus, der $n$ als Eingabe nimmt und die $n$-te Fibonacci-Zahl berechnet. Was beschreibt die Notwendigkeit von Memoization am besten?
\begin{factlist}
    \item b. Memoization ist notwendig, weil Werte (also kleinere Fibonacci-Zahlen) sonst mehrfach berechnet werden müssten.
    \item d. Memoization ist notwendig, weil unsere DP von einer festen, endlichen Anzahl von Werten abhängt.
\end{factlist}
\end{contentbox}


\begin{contentbox}
Gegeben sei folgendes Problem: Gegeben ein Integer-Array, gib die Anzahl der längsten streng monoton steigenden Teilfolgen zurück. Welche der folgenden Aussagen ist FALSCH?
\begin{factlist}
    \item a. Eine effiziente algorithmische Lösung dieses Problems kann nicht als DP kategorisiert werden, weil sie keine Min/Max-Optimierung enthält.
\end{factlist}
\end{contentbox}




\subsection{Wahr oder falsch}\label{sec:13-wahr-oder-falsch}
\begin{contentbox}
Welche der folgenden Aussagen über Basisfälle in einem DP-Algorithmus sind korrekt? Mehrfachauswahl.
\begin{factlist}
    \item a. Basisfälle sind immer notwendig. \textbf{Wahr}
    \item b. Die Korrektheit der durch die Rekurrenz berechneten Werte für die überlappenden Teilprobleme hängt von der Korrektheit der Basisfälle ab. \textbf{Wahr}
    \item c. Basisfälle werden immer einem Wert zugewiesen, ohne Berechnungen oder Vergleiche durchzuführen. \textbf{Falsch}
    \item d. Die Anzahl der Basisfälle muss nie 3 überschreiten. \textbf{Falsch}
\end{factlist}
\end{contentbox}

\begin{contentbox}
Dynamische Programmierung erlaubt es uns immer sofort, einen Teil der Lösung mit einem einzigen Durchlauf der Eingabe zu bestimmen.
\begin{factlist}
    \item b. Falsch
\end{factlist}
\end{contentbox}

\begin{contentbox}
Welche der folgenden Aussagen zur Rückgabe von Lösungen in einem DP-Algorithmus sind immer wahr? Mehrfachauswahl.
\begin{factlist}
    \item a. Sobald eine DP-Tabelle korrekt befüllt ist, genügt ein einziger Lookup, um die optimale Lösung zu bestimmen. \textbf{Wahr}
    \item b. Sobald eine DP-Tabelle korrekt befüllt ist, steht der Wert der optimalen Lösung in einem der Tabelleneinträge oder ist effizient aus den Tabellenwerten berechenbar. \textbf{Wahr}
    \item c. Sobald eine DP-Tabelle korrekt befüllt ist, kann der Wert der optimalen Lösung an mehreren Stellen in der Tabelle stehen. \textbf{Wahr}
\end{factlist}
\end{contentbox}

\begin{contentbox}
Welche der folgenden Anforderungen gelten für die Teilprobleme eines DP-Algorithmus? Mehrfachauswahl.
\begin{factlist}
    \item a. Die Teilprobleme müssen Memoization verwenden, sonst ist keine Lösung möglich. \textbf{Falsch}
    \item b. Die Teilprobleme müssen sich einem Basisfall annähern, sonst ist keine Lösung möglich. \textbf{Wahr}
    \item c. Jedes Teilproblem muss direkt einen Basisfall aufrufen, sonst ist keine Lösung möglich. \textbf{Falsch}
    \item d. Die Teilprobleme müssen alle unterschiedliche Werte für die relevanten Variablen verwenden, sonst ist keine Lösung möglich. \textbf{Falsch}
\end{factlist}
\end{contentbox}

\begin{contentbox}
Angenommen wir haben einen rekursiven DP-Algorithmus, verzichten aber auf Memoization. Welche Probleme können auftreten? Mehrfachauswahl.
\begin{factlist}
    \item a. Der Algorithmus speichert Informationen, die nicht notwendig sind, weil sie später nicht verwendet werden. \textbf{Wahr}
    \item b. Der Algorithmus ist inkorrekt und gibt die falsche Antwort zurück. \textbf{Falsch}
    \item c. Der Algorithmus ist ineffizient und benötigt deutlich mehr Zeit als nötig. \textbf{Wahr}
    \item d. Der Algorithmus löst die Basisfälle nicht korrekt. \textbf{Falsch}
    \item e. Der Algorithmus löst dieselben Teilprobleme mehrfach. \textbf{Wahr}
    \item f. Der Algorithmus verwendet nicht die korrekte Rekurrenz. \textbf{Falsch}
\end{factlist}
\end{contentbox}

\begin{contentbox}
Welche Schritte sind beim Entwurf eines korrekten und effizienten DP-Algorithmus immer notwendig? Mehrfachauswahl.
\begin{factlist}
    \item a. Eine Rekurrenz entwerfen, mit der sich der Wert für alle relevanten Teilprobleme berechnen lässt. \textbf{Wahr}
    \item b. In der Implementierung die Lösungen aller Teilprobleme speichern, die mehr als einmal gelöst werden müssten. \textbf{Wahr}
    \item c. Korrekte Basisfälle für die Rekurrenz haben. \textbf{Wahr}
\end{factlist}
\end{contentbox}

\begin{contentbox}
Für welche Eingabetypen kann dynamische Programmierung ein möglicher Lösungsansatz sein? Mehrfachauswahl.
\begin{factlist}
    \item a. Eingabe: ein Integer-Array. \textbf{Wahr}
    \item b. Eingabe: ein einzelner Integer. \textbf{Wahr}
    \item c. Eingabe: eine zweidimensionale Integer-Matrix. \textbf{Wahr}
    \item d. Eingabe: ein ungerichteter Graph mit $n$ Knoten und $m$ Kanten. \textbf{Wahr}
\end{factlist}
\end{contentbox}

\begin{contentbox}
Dynamische Programmierung liefert für jedes berechenbare Problem eine Lösung in polynomieller Zeit.
\begin{factlist}
    \item b. Falsch
\end{factlist}
\end{contentbox}

\subsection{Code-Fragen}\label{sec:13-code-fragen}

\begin{contentbox}
Frage: Welche der folgenden Implementierungen liefert die $n$-te Tribonacci-Zahl mit der geringsten Zeitkomplexität (asymptotische Gleichheit erlaubt)? Mehrfachauswahl.
\end{contentbox}
\begin{codebox}
\begin{lstlisting}[style=CodeExpert]
function solve(n, seq):
    if seq[n] != NULL:
        return seq[n]
    if n == 0: return 0
    if n == 1 or n == 2: return 1
    seq[n] = solve(n-1, seq) + solve(n-2, seq) + solve(n-3, seq)
    return seq[n]

function main(n):
    1D array seq
    return solve(n, seq)
\end{lstlisting}
\end{codebox}
\begin{codebox}
\begin{lstlisting}[style=CodeExpert]
function main(n):
    if n == 0: return 0
    if n == 1 or n == 2: return 1
    1D array seq
    seq[0] = 0
    seq[1] = 1
    seq[2] = 1
    for i = 3 to n + 1:
        seq[i] = seq[i-1] + seq[i-2] + seq[i-3]
    return seq[n]
\end{lstlisting}
\end{codebox}
\begin{contentbox}
Die $n$-te Fibonacci-Zahl $F(n)$ ist definiert als $F(n) = F(n-1) + F(n-2)$ mit den Basisfällen $F(1) = 1$ und $F(2) = 2$. Dein Kollege schlägt folgende Lösung zur Berechnung der $n$-ten Fibonacci-Zahl vor:
\end{contentbox}
\begin{codebox}
\begin{lstlisting}[style=CodeExpert]
function Fib(n)
    if n == 1: return 1
    else if n == 2: return 1
    else: return Fib(n - 1) + Fib(n - 2)
\end{lstlisting}
\end{codebox}
\begin{contentbox}
Ist die Antwort korrekt und effizient?
\begin{factlist}
    \item b. Die Lösung ist korrekt, aber nicht effizient, weil sie nicht memoisiert und dadurch viele Funktionsaufrufe wiederholt.
\end{factlist}
\end{contentbox}

\begin{contentbox}
Gegeben sei folgendes Problem: Wir möchten die Wahrscheinlichkeit kennen, dass eine bestimmte Anzahl Regentage in den nächsten $n$ Tagen erreicht wird. Für jeden Tag $i$ ist $p_i \in [0, 1]$ die Regenwahrscheinlichkeit. Gegeben sei zudem ein Schwellwert $k$. Gib die Wahrscheinlichkeit zurück, dass es an mindestens $k$ Tagen regnet.
\textbf{Frage:} Du darfst annehmen, dass (a) und (b) beide die korrekte Ausgabe liefern. Welche der beiden ist eine DP-Lösung?
\end{contentbox}
\begin{codebox}
\begin{lstlisting}[style=CodeExpert]
function ChanceOfRain(p[], k)
    2D array memo
    return ChanceOfRainH(p[], 0, k, memo)

function ChanceOfRainH(p[], i, k, memo[][])
    if k == 0:
        return 1
    if i == len(p):
        return 0
    if memo[i][k] != NULL:
        return memo[i][k]
    Rain = p[i] * ChanceOfRainH(p[], i + 1, k - 1, memo)
    Sunny = (1 - p[i]) * ChanceOfRainH(p[], i + 1, k, memo)
    Memo[i][k] = Rain + Sunny
    return Memo[i][k]
\end{lstlisting}
\end{codebox}
\begin{codebox}
\begin{lstlisting}[style=CodeExpert]
function ChanceOfRain(p[], k)
    2D array table
    for i == 0 to k:
        table[n][i] = 0
    for i = n to 1:
        table[i][0] = 1
    for i = n to 1:
        for j = 1 to k:
            rain = table[i+1][j-1] * p[i]
            sunny = table[i+1][j] * (1 - p[i])
            table[i][j] = rain + sunny
    return table[0][k]
\end{lstlisting}
\end{codebox}


\begin{contentbox}
Gegeben sei folgendes Problem: Du möchtest möglichst viel Zeit mit Lesen von Büchern aus deinem Bücherregal verbringen. Du darfst aber kein Buch lesen, das direkt neben einem bereits gelesenen Buch steht. Gegeben ist ein Array mit den Lesezeiten der Bücher. Gib einen Algorithmus an, der die maximale Lesezeit zurückgibt.\\

\textbf{Frage:} Hier ist ein Lösungsversuch:
\end{contentbox}
\begin{codebox}
\begin{lstlisting}[style=CodeExpert]
function BookRead(books[], i, partialTimes[])
    if i >= len(books):
        return 0
    if partialTimes[i] != NULL:
        return partialTimes[i]
    take = BookRead(books, i + 2, partialTimes) + books[i]
    leave = BookRead(books, i + 1, partialTimes)
    return max(take, leave)
\end{lstlisting}
\end{codebox}
\begin{contentbox}
Was verhindert, dass diese Lösung eine DP-Lösung ist?
\begin{factlist}
    \item c. Das Array \texttt{partialTimes} wird nicht korrekt verwendet. \textbf{(korrekt)}
\end{factlist}
\end{contentbox}

\begin{contentbox}
Das Coin-Row-Problem ist wie folgt definiert:
Gegeben ist eine Reihe von $n$ Münzen mit positiven Werten $c_1, c_2, \dots, c_n$. Ziel ist es, Münzen mit maximalem Gesamtwert so auszuwählen, dass keine zwei in der Reihe direkt benachbarten Münzen gemeinsam genommen werden. \\
Sei das Teilproblem D[$i$] die Lösung mit maximalem Wert über die ersten $i$ Münzen: \\
D[$i$] = $\max(c_i + D[i-2], D[i-1])$

\textbf{Beispiel:} \\
Eingabe: [7, 15, 12] $\rightarrow$ Lösung: 7 + 12 = 19 \\
Eingabe: [10, 4, 4, 10, 4, 4, 10] $\rightarrow$ Lösung: 10 + 10 + 10 = 30

\textbf{Lösungsversuch (DP):}
\end{contentbox}
\begin{codebox}
\begin{lstlisting}[style=CodeExpert]
1D Array partialSums
function coinRow(coins, n):
# coins is an array of coin values length n
    if n == 1:
        partialSums[n] = coins[0]
    else if n == 2:
        partialSums[n] = max(coins[0], coins[1])
    else:
        partialSums[n] = max(coins[n] + coinRow(coins, n - 2),
        coinRow(coins, n - 1))
    return partialSums[n]
\end{lstlisting}
\end{codebox}
\begin{contentbox}
Frage: Ist der Code korrekt?
\begin{factlist}
    \item b. Der Code verwendet Memoization nicht korrekt, daher ist die Laufzeit exponentiell.
\end{factlist}
\end{contentbox}


\begin{contentbox}
Gegeben ein positiver Integer $K$. Bestimme die minimale Anzahl Operationen, um $0$ in $K$ zu überführen. Erlaubte Operationen:\\
\begin{factlist}
    \item Eins zum Operanden addieren.
    \item Operanden mit 2 multiplizieren.
\end{factlist}

\textbf{Betrachte diese gültige DP-Lösung:}
\end{contentbox}
\begin{codebox}
\begin{lstlisting}[style=CodeExpert]
function min_operation(k):
    1D Array arr  # array of size k+1
    arr[0] = 0
    for i = 1 TO k:
        arr[i] = arr[i - 1] + 1
        if i % 2 == 0:
            arr[i] = min(arr[i], arr[floor(i / 2)] + 1)
    return arr[k]
\end{lstlisting}
\end{codebox}
\begin{contentbox}
Welche der folgenden Aussagen ist wahr?
\begin{factlist}
    \item c. Die Lösung ist ein gültiger eindimensionaler DP-Ansatz.
\end{factlist}
\end{contentbox}

\begin{contentbox}
Gegeben ist ein Gitter der Grösse $n \times m$, wobei jede Zelle einen nicht-negativen Integer enthält, der die Kosten für das Betreten der Zelle angibt. Startpunkt ist die linke untere Ecke $(n-1, 0)$, Ziel ist die rechte obere Ecke $(0, m-1)$. Minimiere die Gesamtkosten. Lösung:
\end{contentbox}
\begin{codebox}
\begin{lstlisting}[style=CodeExpert]
function minCostPathToTopRight(grid[][]):
    n = len(grid)     # number of rows
    m = len(grid[0])  # number of columns

    # Initialize a 2D DP array to store minimum cost values
    # all elements initialized to 0
    2D Array partialCosts

    # Fill in the DP array
    for i = n - 1 TO 0:
        for j = 0 TO m:
            # Initialize the bottom-left cell
            if i == n - 1 and j == 0:
                partialCosts[i][j] = grid[i][j]
            else if i == n - 1:
                partialCosts[i][j] = grid[i][j] +
                partialCosts[i][j - 1]
            else if j == 0:
                partialCosts[i][j] = grid[i][j] +
                partialCosts[i + 1][j]
            else:
                partialCosts[i][j] = grid[i][j] +
                min(partialCosts[i][j - 1],
                partialCosts[i + 1][j])

    return partialCosts[0][m - 1]
\end{lstlisting}
\end{codebox}
\begin{contentbox}
Wähle die korrekte Aussage:
\begin{factlist}
    \item b. Dies ist eine gültige DP-Lösung, die das Gitter von links unten nach rechts oben durchläuft.
\end{factlist}
\end{contentbox}

\begin{contentbox}
Das Stock-Market-Problem ist wie folgt definiert: Gegeben eine Sequenz historischer Tageskurse einer Aktie über $n$ aufeinanderfolgende Tage. Bestimme die Tage, an denen man die Aktie hätte kaufen und verkaufen sollen, um den Gewinn zu maximieren, und gib den maximalen Gewinn zurück.
Hinweis: Es darf nur einmal gekauft und einmal verkauft werden, und der Kauftag muss vor oder gleich dem Verkaufstag liegen. \\
Beispiele: \\
Eingabe: [2, 5, 7, 1] $\rightarrow$ Lösung: 5 (Kauf an Tag 1, Verkauf an Tag 3)\\
Eingabe: [5, 2, 4, 8, 5, 6, 7] $\rightarrow$ Lösung: 6 (Kauf an Tag 2, Verkauf an Tag 4)\\
\textbf{Welche der folgenden Funktionen ist eine DP-Lösung des Stock-Market-Problems?}
\end{contentbox}
\begin{codebox}
\begin{lstlisting}[style=CodeExpert]
function maxProfit(prices[]):
    minPrice = min(prices)
    minIndex = len(prices)
    maxPrice = 0
    for i = 1 to len(prices):
        if prices[i] == minPrice:
            minIndex = i
        if minIndex < i and prices[i] > maxPrice:
            maxPrice = prices[i]
    return max(0, maxPrice - minPrice)
\end{lstlisting}
\end{codebox}
\begin{codebox}
\begin{lstlisting}[style=CodeExpert]
function maxProfit(prices[]):
    if (len(prices) == 0):
        return 0
    maxProfit = 0
    minPrice = prices[0]
    for i = 1 to len(prices):
        if prices[i] < minPrice:
            minPrice = prices[i]
        else:
            maxProfit = max(maxProfit, prices[i] - minPrice)
    return maxProfit
\end{lstlisting}
\end{codebox}

\begin{contentbox}
Das Longest-Common-Subsequence-Problem (LCS) ist wie folgt definiert: Gegeben zwei Strings (Zeichen-Arrays) X[0..m-1] und Y[0..n-1] der Länge $m$ und $n$. Finde die Länge der längsten Teilfolge, die in X und Y vorkommt.
Beispiel: \\
X = [a, b, c, d, e, f] \\
Y = [d, z, a, f, c, w, d, e] \\
Antwort = 4 (Teilfolge [a, c, d, e])\\
\textbf{Welche der folgenden Funktionen löst das LCS-Problem mit der geringsten Zeitkomplexität? Eine Antwort.}
\end{contentbox}
\begin{codebox}
\begin{lstlisting}[style=CodeExpert]
# arr is of appropriate size
function lcs(X, Y, m, n, arr[]):
    if (m == 0 or n == 0):
        return 0
    if (arr[m][n] != NULL):
        return arr[m][n]
    if X[m - 1] == Y[n - 1]:
        arr[m][n] = 1 + lcs(X, Y, m - 1, n - 1, arr)
        return arr[m][n]
    arr[m][n] = max(lcs(X, Y, m, n - 1, arr), lcs(X, Y,
    m - 1, n, arr))
    return arr[m][n]
\end{lstlisting}
\end{codebox}
\begin{codebox}
\begin{lstlisting}[style=CodeExpert]
# arr is of appropriate size
function lcs(X, Y, m, n, arr[]):
    for i = 0 to m:
        for j = 0 to n:
            if i == 0 or j == 0:
                arr[i][j] = 0
            else if X[i - 1] == Y[j - 1]:
                arr[i][j] = arr[i - 1][j - 1] + 1
            else:
                arr[i][j] = max(arr[i - 1][j], arr[i][j - 1])
    return arr[m][n]
\end{lstlisting}
\end{codebox}



\begin{codebox}
\begin{lstlisting}[style=CodeExpert, extendedchars=true, inputencoding=utf8]
# print(dir(math)) falls syntax unbekannt
import math
c=0
def f(): global c; c+=1
def check(g):
    S=[10,14,18,22]; C=[]
    for n in S: global c; c=0; g(n); C.append(c)
    T={"log n":math.log2, "sqrt(n)":lambda n:n**.5}
    E={}
    for k,h in T.items():
        try:
            R=[C[i]/h(n) for i,n in enumerate(S) if h(n)>1e-9]
            if not R: continue
            M=sum(R)/len(R)
            if M<1e-9: continue
            E[k]=math.sqrt(sum((x-M)**2 for x in R)/len(R))/M
        except: continue
    return "theta("+min(E,key=E.get)+")" if E else "theta(?)"

def g(n): # g von Aufgabe einfuegen
    pass

print(check(g)) # Ergebnis printen
\end{lstlisting}
\end{codebox}



